% Кодировка и язык
\usepackage[utf8]{inputenc}           % Кодировка исходного файла (UTF-8)
\usepackage[T2A]{fontenc}             % Кодировка шрифтов для кириллицы
\usepackage[english,russian]{babel}   % Локализация: русские названия, переносы

% Поиск в PDF
\usepackage{cmap}

% Поля страницы
\usepackage{geometry}
\geometry{left=3.8cm, right=3.8cm, top=2.5cm, bottom=2.5cm, paperheight=29.7cm}

% Математика
\usepackage{amsmath, amssymb, amsthm} % Базовые пакеты для математики
\usepackage{mathtools}                % Расширение amsmath (доп. символы и окружения)

% Картинки
\usepackage{graphicx}                 % Вставка картинок (\includegraphics)
\graphicspath{{pictures/}}            % Путь к папке с картинками
\usepackage{float}                    % Позволяет использовать [H] для жёсткого позиционирования
\usepackage{wrapfig}                  % Обтекание текста вокруг картинок

% Таблицы
\usepackage{array}                    % Расширенные возможности для таблиц
\usepackage{booktabs}                 % Красивые таблицы (\toprule, \midrule, \bottomrule)
\usepackage{multirow}                 % Объединение строк в таблицах
\usepackage{tabularx}                 % Таблицы с автоматической шириной колонок
\usepackage{ragged2e}                 % Улучшенное выравнивание (используется в tabularx)

% Списки
\usepackage{enumitem}                 % Гибкая настройка списков (itemize, enumerate)

% Колонтитулы
\usepackage{fancyhdr}

% Гиперссылки
\usepackage{xcolor}                   % Работа с цветами
\usepackage{hyperref}                 % Гиперссылки, оглавление, перекрёстные ссылки

% Умные ссылки
\usepackage{cleveref}

% Типографика
\usepackage{microtype}                % Улучшает переносы и межсловные расстояния

% Подписи к картинкам и таблицам
\usepackage{caption}                  % Настройка подписей (шрифт, отступы)

% Рисование схем и графиков
\usepackage{tikz}                     % Рисование схем, графиков, диаграмм
\usepackage{pgfplots}                 % Построение графиков функций (на основе tikz)
\pgfplotsset{compat=newest}           % Версия совместимости pgfplots

% Алгоритмы и код
\usepackage{algorithm}                % Окружения для алгоритмов
\usepackage{algorithmic}              % Синтаксис для алгоритмов
\usepackage{listings}                 % Вставка кода с подсветкой синтаксиса

\usepackage{soul}                     % Выделение текста: подчёркивание, зачёркивание, подсветка
\usepackage{multicol}                 % Многоколоночная вёрстка текста
\usepackage{titlesec}                 % Настройка внешнего вида заголовков секций
\usepackage{parskip}                  % Отступы между абзацами вместо красной строки
\usepackage{setspace}                 % Настройка межстрочного интервала
\usepackage{tcolorbox}                % Цветные блоки для заметок, цитат и выделений

% Палитра Quiet Light из VS Code
\definecolor{quietLightForeground}{HTML}{333333}
\definecolor{quietLightEditorBg}{HTML}{F5F5F5}
\definecolor{quietLightPanelBg}{HTML}{F2F2F2}
\definecolor{quietLightBorder}{HTML}{C9D0D9}
\definecolor{quietLightComment}{HTML}{AAAAAA}
\definecolor{quietLightKeyword}{HTML}{4B69C6}
\definecolor{quietLightType}{HTML}{7A3E9D}
\definecolor{quietLightConstant}{HTML}{9C5D27}
\definecolor{quietLightString}{HTML}{448C27}
\definecolor{quietLightFunction}{HTML}{AA3731}
\definecolor{quietLightPunctuation}{HTML}{777777}

% Совместимость со старыми именами цветов
\colorlet{pblue}{quietLightKeyword}
\colorlet{pgreen}{quietLightString}
\colorlet{pred}{quietLightFunction}
\colorlet{pgrey}{quietLightPunctuation}

\onehalfspacing

\titleformat{\section}
  {\Large\bfseries}
  {\thesection}
  {0.7em}
  {}

\titleformat{\subsection}
  {\large\bfseries}
  {\thesubsection}
  {0.7em}
  {}

\lstset{
    basicstyle=\ttfamily\small\color{quietLightForeground},
    backgroundcolor=\color{quietLightEditorBg},
    frame=single,
    rulecolor=\color{quietLightBorder},
    breaklines=true,
    columns=fullflexible,
    keepspaces=true,
    showstringspaces=false,
    keywordstyle=\color{quietLightKeyword}\bfseries,
    morekeywords=[2]{define,elif,else,endif,error,if,ifdef,ifndef,include,line,pragma,undef,warning},
    keywordstyle=[2]\color{quietLightKeyword},
    deletekeywords={bool,char,double,float,int,long,short,signed,unsigned,void},
    morekeywords=[3]{bool,char,double,float,int,long,short,signed,unsigned,void,size_t,string},
    keywordstyle=[3]\color{quietLightType},
    identifierstyle=\color{quietLightType},
    commentstyle=\color{quietLightComment}\itshape,
    stringstyle=\color{quietLightString},
    numberstyle=\color{quietLightConstant},
}

\lstdefinestyle{nolint}{
    basicstyle=\ttfamily\small\color{quietLightForeground},
    backgroundcolor=\color{quietLightPanelBg},
    columns=fullflexible,
    keepspaces=true,
    breaklines=false,
    keywordstyle=\color{quietLightKeyword},
    commentstyle=\color{quietLightKeyword},
    stringstyle=\color{quietLightKeyword},
    numberstyle=\color{quietLightKeyword},
}

\newcommand{\nolint}{\lstinline[style=nolint]}

\newtcolorbox{habrquote}{
    colback=quietLightPanelBg,
    colframe=quietLightBorder,
    boxrule=0.4pt,
    arc=2pt,
    left=8pt,
    right=8pt,
    top=6pt,
    bottom=6pt
}

% Настройка нумерации секций и подсекций
\renewcommand{\thesection}{\arabic{section}}
\renewcommand{\thesubsection}{\thesection.\arabic{subsection}}


% Настройка списков: убираем лишние отступы слева
\setlist[itemize]{leftmargin=*}
\setlist[enumerate]{leftmargin=*}


% Определение окружений для теорем, лемм, определений и т.д.
% Нумерация привязана к номеру секции
\newtheorem{theorem}{Теорема}[section]
\newtheorem{definition}{Определение}[section]
\newtheorem{remark}{Замечание}[section]
\newtheorem{lemma}{Лемма}[section]
\newtheorem{example}{Пример}[section]
\newtheorem{consequence}{Следствие}[section]
\newtheorem{algorithmm}{Алгоритм}
\newtheorem{question}{Вопрос}[section]
\newtheorem{answer}{Ответ}[section]
\newtheorem{statement}{Утверждение}[section]
